\honest{The Level Above Mine}{Eliezer Yudkowsky}{2008}

A short detour. I lent Mike Li my copy of Jaynes's \textsc{Probability Theory: The Logic of Science}, and he came back saying: ``Wow\ldots\ it's like Jaynes is a thousand-year-old vampire.'' I knew what he meant: in fantasy, the older a vampire gets, the more powerful. I had enjoyed proofs before, but Jaynes was the first to give me a sense of ``formidability'' from mathematics, maybe because he lined up the paradoxes used against Bayesianism and blasted them to pieces, or because he cared about probability theory instead of treating it as a game of aesthetics. He seemed to reach the right answer by the shortest route, tearing the surrounding mistakes to shreds in the same motion. \nb{The book has a chapter on them, ``Paradoxes of probability theory''. Not all stayed in pieces: the authors of the marginalization paradox published a critique of Jaynes's treatment in 1996, and Timothy Wallstrom (2003), who agrees with Jaynes at the conceptual level, found his technical analysis ``not entirely successful''.} Of course, a book shows only its author's best side. \nb{Jaynes did not finish this one. It was published in 2003, five years after his death, from an incomplete manuscript edited by Larry Bretthorst.}

It spoke well of Mike that he sensed this. ``It's a general rule, I've observed, that you can't discriminate between levels too far above your own.'' Someone once earnestly told me that I was really bright and ``ought to go to college''. The limit may be about one standard deviation, ``though that's just a cool-sounding wild guess.'' \nb{Kruger and Dunning (1999) found that the weakest performers on their tests were ``less able to recognize competence in others'' than the strongest, though they compared only those two groups.}

I asked Mike whether he got the same sense from me. He said no. I had not thought I had reached Jaynes's level. I aspire to it, in AI and reflectivity. I could plead that my art is harder than his, but ``there is no art of which I am as much the master now, as Jaynes was of probability theory.'' This does not place me beneath Jaynes as a person. It recognizes a level of expertise I have not reached: I can argue forcefully in my subject, but that is not writing out the equations and saying ``DONE.'' Until I reach it, I must allow that my native talent may never be enough.

I asked Marcello Herreshoff whether anyone struck him as more natively intelligent than me. He named John Conway, whom he had met at a summer math camp: ``He just struck me as having a tremendous amount of mental horsepower.'' Worse, it was an ultra-famous older man I could not recruit. Perhaps I have not had the chance to show a mastered subject; perhaps our minds point different ways; or perhaps ``I'm strictly dumber than Conway, dominated by him along all dimensions.'' That hurts to confess. I have accepted it, because I have acted on it: I blog partly so that a younger mind can read this and ``zip off'' past me. If I wasted my youth setting up support for myself instead of studying math full time, I may hit a brick wall at forty, with nothing left but to pass the resources on to a younger mind. So I should at least leave that mind a trail to my successes and warnings on my mistakes.

Such ``specific efforts'' are the only humility I credit myself with, along with giving up my precious theories when they failed the standard Jaynes had shown me; that was hard, and real. ``Modest demeanors are cheap.'' Too many people, faced with a counterargument, say ``I am but a fallible mortal, of course I could be wrong'' and then do exactly what they planned. \nb{This is the definition the author gave in ``The Proper Use of Humility'' (2006): ``To be humble is to take specific actions in anticipation of your own errors.''}

I do not say that someone less brilliant than Jaynes or Conway can still do important things in my field. ``Because I do know\ldots\ that's not how it works.'' I do not explain here how it works. \nb{The explanation came nine days later, in ``My Bayesian Enlightenment'': Friendly AI, as the author saw it, needed ``something like the Jaynes-level'' of mastery.}
